Optimised product quantisation (OPQ) learns a rotation and subcodebooks
to reduce reconstruction error \cite{opq}. ScaNN's anisotropic
quantisation instead targets inner-product approximation \cite{scann}.
The distinction matters because minimising vector reconstruction error
need not best preserve neighbour ordering.

RaBitQ estimates distances using randomised binary codes \cite{rabitq},
with a multi-bit extension offering a finer representation \cite{rabitq_multi}.
GPU IVF-RaBitQ combines compressed distance estimation with IVF routing
and fused search \cite{ivfrabitq}. Both the original one-bit RaBitQ
construction and JHQ's evaluated primary code store signs after rotation.
JHQ uses fixed reconstruction levels and residual-code refinement;
the IVF-RaBitQ baseline tests 4/8/9 bits per dimension.
\Cref{sec:7-2,sec:7-smallbatch} compare the complete systems'
recall--throughput trade-offs.

Fast Scan and Quicker ADC accelerate CPU distance computation using
small tables accessed through single-instruction, multiple-data (SIMD)
operations \cite{fastscan,quickeradc}. Table size and hardware lookup support are therefore central to speed.
JHQ-GPU splits tables using its primary codebook's coordinate structure;
Chapter 5 analyses the identity and finite-precision effects.

Faiss and the GPU search library cuVS expose parameters controlling the
number of candidates refined \cite{faiss_refine,cuvs_refinement}.
Faiss autotuning uses labelled evaluation \cite{faiss_autotune}; learned
early termination predicts when to stop \cite{li_termination}; DARTH
expresses search effort through recall targets \cite{darth}.
These approaches provide the comparison points for the label-free
calibration in Chapter 6.
